\documentclass[12pt]{article}
\usepackage[utf8]{inputenc}
\usepackage[T1]{fontenc}
\usepackage{hyperref}
\title{Phala Language Resources for K-12 Curriculum Development}
\author{Compiled by Claude Sonnet 5.5 (claude-sonnet-5-5) for Andr\'as Kornai}
\date{}
\begin{document}
\maketitle

\section{Introduction}
Phala (ISO 639-3 \texttt{ypa}; Glottocode phal1256) is a Loloish (Sino-Tibetan, Lolo-Burmese) language of southwestern China, in the Riverine Phula group of Southeastern Ngwi. Sources give about 12,000 to 13,000 speakers in Honghe, Shiping, Jianshui and Yuanyang Counties, Yunnan. The population figures vary between sources. The dialect documented by Pelkey (2011) is from Luodie village, Yuanjiang County. Phala should not be confused with Pholo (\texttt{yip}), a separate language.

\section{Primary Educational Resources}
No primers, textbooks, or school materials were found in search results.

\section{Linguistic and Reference Materials}
\begin{itemize}
\item Wikipedia, ``Phola language'' (cluster including Phala): \url{https://en.wikipedia.org/wiki/Phola_language}
\item Wikipedia, ``Southeastern Loloish languages'': \url{https://en.wikipedia.org/wiki/Southeastern_Loloish_languages}
\item Wikipedia, ``Loloish languages'': \url{https://en.wikipedia.org/wiki/Loloish_languages}
\item Glottolog entry for Phala: \url{https://glottolog.org/resource/languoid/id/phal1256.bigmap.html}
\end{itemize}

\section{Religious and Community Resources}
None found.

\section{Digital and Audio Resources}
None found.

\section{Curriculum Development Context}
Phala is documented mainly in comparative linguistic work, notably Pelkey (2011), which the search summary cited but I did not locate directly. No educational material exists in the results, and nothing found indicates school use. A curriculum would need new documentation and community collaboration.
\end{document}